% \clearpage
\section*{Discussion: Intrastate Conflict as a Dynamic Network}

Intrastate conflicts are often more complex than Manichaean struggles between the government and a unified opposition. More and more, they involve a number of actors operating in an environment where interactions are interdependent. The improved predictive performance of our network approach is a result of the fact that there is significant network based structure underlying the Nigerian conflict system. Importantly, this network structure is more than just a useful part of our modeling approach; it actually helps to explain who fights who during conflict. 

In our study, first order dependencies reveal that ethnic-militias in central Nigeria are more violent than expected. Interestingly, one of the striking findings is that the Fulani militia is particularly more active in engaging in battles than would be predicted by their relationships with the civilian population, countrywide factors, and their indirect links to other groups. The first order effects also reveal that the center and southwest regions of the country feature relatively less struggles with the government, compared to other areas, and more intra-rebel group violence. These first order effects may be because, in ideologically diverse regions, groups operate in a target-rich environment, where there are more opportunities to advance their ideology through violence. Second order effects, such as reciprocity also reveal important insights. The presence of reciprocity in battles, and the consistent relationship between civilian victimization and conflict, confirm our expectation that the interactions between actors matter at least as much as their mere characteristics.

Because we are able to estimate a network model in which actor composition changes over time, we can assess the impact that particular actors such as Boko Haram have on the system. We find that the entrance of this single actor profoundly influenced the occurrence of battles across the system. Boko Haram's entrance does not simply increase conflict \emph{directly}, it also is associated with a marked rise in violence in the dyads that \emph{do not include Boko Haram.} This raises an important question for future research which is whether or not we can formally study why some actors have stronger effects on shaping interactions across a network than others. 

\section*{Conclusion}
These contributions are not without their limitations. While our study affirms the utility of ACLED data efforts, more actor-level characteristics are needed to fully engage the rich literature on civil war dynamics. Such data would enable even more precise investigation of covariates of interest for determining changes in rebel group behavior over time. Second, though our study has revealed meaningful findings with respect to key covariates such as civilian victimization, more research is needed to explore the generalizability of such results. Even after controlling for network related dependencies within this system, we find that actors who target civilians are more likely to receive and send conflict themselves. Relatedly, we find little evidence that civilian led protest against a particular actor changes that actor's likelihood in initiating conflict.  The findings with regards to the role of civilian victimization and civilian riots/protests might not generalize beyond the Nigerian case. Nigeria may be unique for several reasons, its 150 million people are divided almost equally between Muslims and Christians. Further, Nigerians split into even finer divisions based on tribe, of which there are over 200 in the country. Future research should compare and contrast these findings beyond the Nigerian case to better illuminate the networked nature of civilian mobilization in conflict environments.  We hope this study motivates future work that more fully considers the relationship between the behavior of armed actors and civilian behavior in complex settings.

Our study has embraced the complexity of multi-actor conflicts, rather than discounted it. We have shown that network dynamics help explain the occurrence of violence between groups over time, even as actors enter and exit the conflict. We have demonstrated how key, violent actors can increase violence within a network system. Additionally, our approach achieves a longstanding goal of conflict studies: to accurately predict violence over time. Our approach offers important insights into how conflict studies can utilize methodological innovations to expand theoretical inquiry. 
